% Source: book pp. 163--174 (PDF 177--188); no solutions (manual stops at 5B).
% \axtag and \exnote: see the top of 5C.tex.
\DeclareRobustCommand\axtag[1]{\refstepcounter{theorem}\tag{\thetheorem}\label{#1}}
\section{Diagonalizable Operators}

\subsection{Diagonal Matrices}

\begin{definition}[Diagonal matrix]\label{ax:5.48}
A \emph{diagonal matrix} is a square matrix that is $0$ everywhere except
possibly on the diagonal.
\end{definition}

\begin{example}[Diagonal matrix]\label{ax:5.49}\leavevmode
\[ \begin{pmatrix} 8&0&0\\ 0&5&0\\ 0&0&5 \end{pmatrix} \]
is a diagonal matrix.
\end{example}

\marginnote{Every diagonal matrix is upper triangular. Diagonal matrices
typically have many more $0$'s than most upper-triangular matrices of the same
size.}%
If an operator has a diagonal matrix with respect to some basis, then the entries
on the diagonal are precisely the eigenvalues of the operator; this follows from
\ref{ax:5.41} (or find an easier direct proof for diagonal matrices).

\begin{definition}[Diagonalizable]\label{ax:5.50}
An operator on $V$ is called \emph{diagonalizable} if the operator has a diagonal
matrix with respect to some basis of $V$.
\end{definition}

\begin{example}[Diagonalization may require a different basis]\label{ax:5.51}
Define $T\in\Lin(\R^2)$ by
\[ T(x,y)=(41x+7y,-20x+74y). \]
The matrix of $T$ with respect to the standard basis of $\R^2$ is
\[ \begin{pmatrix} 41&7\\ -20&74 \end{pmatrix}, \]
which is not a diagonal matrix. However, $T$ is diagonalizable. Specifically, the
matrix of $T$ with respect to the basis $(1,4)$, $(7,5)$ is
\[ \begin{pmatrix} 69&0\\ 0&46 \end{pmatrix} \]
because $T(1,4)=(69,276)=69(1,4)$ and $T(7,5)=(322,230)=46(7,5)$.
\end{example}

For $\lambda\in\F$, we will find it convenient to have a name and a notation for
the set of vectors that an operator $T$ maps to $\lambda$ times the vector.

\begin{definition}[Eigenspace, $E(\lambda,T)$]\label{ax:5.52}
Suppose $T\in\Lin(V)$ and $\lambda\in\F$. The \emph{eigenspace} of $T$
corresponding to $\lambda$ is the subspace $E(\lambda,T)$ of $V$ defined by
\[ E(\lambda,T)=\nul(T-\lambda I)=\{v\in V : Tv=\lambda v\}. \]
Hence $E(\lambda,T)$ is the set of all eigenvectors of $T$ corresponding to
$\lambda$, along with the $0$ vector.
\end{definition}

For $T\in\Lin(V)$ and $\lambda\in\F$, the set $E(\lambda,T)$ is a subspace of
$V$ because the null space of each linear map on $V$ is a subspace of $V$. The
definitions imply that $\lambda$ is an eigenvalue of $T$ if and only if
$E(\lambda,T)\ne\{0\}$.

\begin{example}[Eigenspaces of an operator]\label{ax:5.53}
Suppose the matrix of an operator $T\in\Lin(V)$ with respect to a basis
$v_1,v_2,v_3$ of $V$ is the matrix in Example~\ref{ax:5.49}. Then
\[ E(8,T)=\spn(v_1), \qquad E(5,T)=\spn(v_2,v_3). \]
\end{example}

If $\lambda$ is an eigenvalue of an operator $T\in\Lin(V)$, then $T$ restricted
to $E(\lambda,T)$ is just the operator of multiplication by $\lambda$.

\begin{result}[Sum of eigenspaces is a direct sum]\label{ax:5.54}
Suppose $T\in\Lin(V)$ and $\lambda_1,\dots,\lambda_m$ are distinct eigenvalues
of $T$. Then
\[ E(\lambda_1,T)+\cdots+E(\lambda_m,T) \]
is a direct sum. Furthermore, if $V$ is finite-dimensional, then
\[ \dim E(\lambda_1,T)+\cdots+\dim E(\lambda_m,T)\le\dim V. \]
\end{result}

\begin{proof}
To show that $E(\lambda_1,T)+\cdots+E(\lambda_m,T)$ is a direct sum, suppose
\[ v_1+\cdots+v_m=0, \]
where each $v_k$ is in $E(\lambda_k,T)$. Because eigenvectors corresponding to
distinct eigenvalues are linearly independent (by \ref{ax:5.11}), this implies
that each $v_k$ equals $0$. Thus $E(\lambda_1,T)+\cdots+E(\lambda_m,T)$ is a
direct sum (by \ref{ax:1.45}), as desired.

Now suppose $V$ is finite-dimensional. Then
\begin{align*}
\dim E(\lambda_1,T)+\cdots+\dim E(\lambda_m,T)
  &= \dim\bigl(E(\lambda_1,T)\oplus\cdots\oplus E(\lambda_m,T)\bigr)\\
  &\le \dim V,
\end{align*}
where the first line follows from \ref{ax:3.94} and the second line follows from
\ref{ax:2.37}.
\end{proof}

\subsection{Conditions for Diagonalizability}

The following characterizations of diagonalizable operators will be useful.

\begin{result}[Conditions equivalent to diagonalizability]\label{ax:5.55}
Suppose $V$ is finite-dimensional and $T\in\Lin(V)$. Let
$\lambda_1,\dots,\lambda_m$ denote the distinct eigenvalues of $T$. Then the
following are equivalent.
\begin{parts}
\item $T$ is diagonalizable.
\item $V$ has a basis consisting of eigenvectors of $T$.
\item $V=E(\lambda_1,T)\oplus\cdots\oplus E(\lambda_m,T)$.
\item $\dim V=\dim E(\lambda_1,T)+\cdots+\dim E(\lambda_m,T)$.
\end{parts}
\end{result}

\begin{proof}
An operator $T\in\Lin(V)$ has a diagonal matrix
\[ \begin{pmatrix} \lambda_1 & & 0\\ & \ddots & \\ 0 & & \lambda_n
   \end{pmatrix} \]
with respect to a basis $v_1,\dots,v_n$ of $V$ if and only if
$Tv_k=\lambda_kv_k$ for each $k$. Thus (a) and (b) are equivalent.

Suppose (b) holds; thus $V$ has a basis consisting of eigenvectors of $T$. Hence
every vector in $V$ is a linear combination of eigenvectors of $T$, which implies
that
\[ V=E(\lambda_1,T)+\cdots+E(\lambda_m,T). \]
Now \ref{ax:5.54} shows that (c) holds, proving that (b) implies (c).

That (c) implies (d) follows immediately from \ref{ax:3.94}.

Finally, suppose (d) holds; thus
\[ \dim V=\dim E(\lambda_1,T)+\cdots+\dim E(\lambda_m,T). \axtag{ax:5.56} \]
Choose a basis of each $E(\lambda_k,T)$; put all these bases together to form a
list $v_1,\dots,v_n$ of eigenvectors of $T$, where $n=\dim V$ (by
\ref{ax:5.56}). To show that this list is linearly independent, suppose
\[ a_1v_1+\cdots+a_nv_n=0, \]
where $a_1,\dots,a_n\in\F$. For each $k=1,\dots,m$, let $u_k$ denote the sum of
all the terms $a_jv_j$ such that $v_j\in E(\lambda_k,T)$. Thus each $u_k$ is in
$E(\lambda_k,T)$, and
\[ u_1+\cdots+u_m=0. \]
Because eigenvectors corresponding to distinct eigenvalues are linearly
independent (see \ref{ax:5.11}), this implies that each $u_k$ equals $0$.
Because each $u_k$ is a sum of terms $a_jv_j$, where the $v_j$'s were chosen to
be a basis of $E(\lambda_k,T)$, this implies that all $a_j$'s equal $0$. Thus
$v_1,\dots,v_n$ is linearly independent and hence is a basis of $V$ (by
\ref{ax:2.38}). Thus (d) implies (b), completing the proof.
\end{proof}

For additional conditions equivalent to diagonalizability, see \ref{ax:5.62},
Exercises~5 and~15 in this section, Exercise~24 in Section~7B, and Exercise~15
in Section~8A.

As we know, every operator on a nonzero finite-dimensional complex vector space
has an eigenvalue. However, not every operator on a nonzero finite-dimensional
complex vector space has enough eigenvectors to be diagonalizable, as shown by
the next example.

\begin{example}[An operator that is not diagonalizable]\label{ax:5.57}
Define an operator $T\in\Lin(\F^3)$ by $T(a,b,c)=(b,c,0)$. The matrix of $T$
with respect to the standard basis of $\F^3$ is
\[ \begin{pmatrix} 0&1&0\\ 0&0&1\\ 0&0&0 \end{pmatrix}, \]
which is an upper-triangular matrix but is not a diagonal matrix.

As you should verify, $0$ is the only eigenvalue of $T$ and furthermore
\[ E(0,T)=\{(a,0,0)\in\F^3 : a\in\F\}. \]
Hence conditions (b), (c), and (d) of \ref{ax:5.55} fail (of course, because
these conditions are equivalent, it is sufficient to check that only one of them
fails). Thus condition (a) of \ref{ax:5.55} also fails. Hence $T$ is not
diagonalizable, regardless of whether $\F=\R$ or $\F=\C$.
\end{example}

The next result shows that if an operator has as many distinct eigenvalues as the
dimension of its domain, then the operator is diagonalizable.

\begin{result}[Enough eigenvalues implies diagonalizability]\label{ax:5.58}
Suppose $V$ is finite-dimensional and $T\in\Lin(V)$ has $\dim V$ distinct
eigenvalues. Then $T$ is diagonalizable.
\end{result}

\begin{proof}
Suppose $T$ has distinct eigenvalues $\lambda_1,\dots,\lambda_{\dim V}$. For
each $k$, let $v_k\in V$ be an eigenvector corresponding to the eigenvalue
$\lambda_k$. Because eigenvectors corresponding to distinct eigenvalues are
linearly independent (see \ref{ax:5.11}), $v_1,\dots,v_{\dim V}$ is linearly
independent.

A linearly independent list of $\dim V$ vectors in $V$ is a basis of $V$ (see
\ref{ax:2.38}); thus $v_1,\dots,v_{\dim V}$ is a basis of $V$. With respect to
this basis consisting of eigenvectors, $T$ has a diagonal matrix.
\end{proof}

In later chapters we will find additional conditions that imply that certain
operators are diagonalizable. For example, see the real spectral theorem
(\ref{ax:7.29}) and the complex spectral theorem (\ref{ax:7.31}).

The result above gives a sufficient condition for an operator to be
diagonalizable. However, this condition is not necessary. For example, the
operator $T$ on $\F^3$ defined by $T(x,y,z)=(6x,6y,7z)$ has only two eigenvalues
($6$ and $7$) and $\dim\F^3=3$, but $T$ is diagonalizable (by the standard basis
of $\F^3$).

\marginnote{For a spectacular application of these techniques, see Exercise~21,
which shows how to use diagonalization to find an exact formula for the
$n^{\text{th}}$ term of the Fibonacci sequence.}%
The next example illustrates the importance of diagonalization, which can be used
to compute high powers of an operator, taking advantage of the equation
$T^kv=\lambda^kv$ if $v$ is an eigenvector of $T$ with eigenvalue $\lambda$.

\begin{example}[Using diagonalization to compute $T^{100}$]\label{ax:5.59}
Define $T\in\Lin(\F^3)$ by $T(x,y,z)=(2x+y,5y+3z,8z)$. With respect to the
standard basis, the matrix of $T$ is
\[ \begin{pmatrix} 2&1&0\\ 0&5&3\\ 0&0&8 \end{pmatrix}. \]
The matrix above is an upper-triangular matrix but it is not a diagonal matrix.
By \ref{ax:5.41}, the eigenvalues of $T$ are $2$, $5$, and $8$. Because $T$ is an
operator on a vector space of dimension three and $T$ has three distinct
eigenvalues, \ref{ax:5.58} assures us that there exists a basis of $\F^3$ with
respect to which $T$ has a diagonal matrix.

To find this basis, we only have to find an eigenvector for each eigenvalue. In
other words, we have to find a nonzero solution to the equation
\[ T(x,y,z)=\lambda(x,y,z) \]
for $\lambda=2$, then for $\lambda=5$, and then for $\lambda=8$. Solving these
simple equations shows that for $\lambda=2$ we have an eigenvector $(1,0,0)$, for
$\lambda=5$ we have an eigenvector $(1,3,0)$, and for $\lambda=8$ we have an
eigenvector $(1,6,6)$.

Thus $(1,0,0)$, $(1,3,0)$, $(1,6,6)$ is a basis of $\F^3$ consisting of
eigenvectors of $T$, and with respect to this basis the matrix of $T$ is the
diagonal matrix
\[ \begin{pmatrix} 2&0&0\\ 0&5&0\\ 0&0&8 \end{pmatrix}. \]

To compute $T^{100}(0,0,1)$, for example, write $(0,0,1)$ as a linear
combination of our basis of eigenvectors:
\[ (0,0,1)=(1/6)(1,0,0)-(1/3)(1,3,0)+(1/6)(1,6,6). \]
Now apply $T^{100}$ to both sides of the equation above, getting
\begin{align*}
T^{100}(0,0,1)
  &= 1/6\bigl(T^{100}(1,0,0)\bigr)-1/3\bigl(T^{100}(1,3,0)\bigr)
     +1/6\bigl(T^{100}(1,6,6)\bigr)\\
  &= 1/6\bigl(2^{100}(1,0,0)-2\cdot5^{100}(1,3,0)+8^{100}(1,6,6)\bigr)\\
  &= 1/6\bigl(2^{100}-2\cdot5^{100}+8^{100},\
     6\cdot8^{100}-6\cdot5^{100},\ 6\cdot8^{100}\bigr).
\end{align*}
\end{example}

We saw earlier that an operator $T$ on a finite-dimensional vector space $V$ has
an upper-triangular matrix with respect to some basis of $V$ if and only if the
minimal polynomial of $T$ equals $(z-\lambda_1)\cdots(z-\lambda_m)$ for some
$\lambda_1,\dots,\lambda_m\in\F$ (see \ref{ax:5.44}). As we previously noted
(see \ref{ax:5.47}), this condition is always satisfied if $\F=\C$.

Our next result \ref{ax:5.62} states that an operator $T\in\Lin(V)$ has a
diagonal matrix with respect to some basis of $V$ if and only if the minimal
polynomial of $T$ equals $(z-\lambda_1)\cdots(z-\lambda_m)$ for some
\textbf{\emph{distinct}} $\lambda_1,\dots,\lambda_m\in\F$. Before formally
stating this result, we give two examples of using it.

\begin{example}[Diagonalizable, but with no known exact eigenvalues]%
\label{ax:5.60}
Define $T\in\Lin(\C^5)$ by
\[ T(z_1,z_2,z_3,z_4,z_5)=(-3z_5,z_1+6z_5,z_2,z_3,z_4). \]
The matrix of $T$ is shown in Example~\ref{ax:5.26}, where we showed that the
minimal polynomial of $T$ is $3-6z+z^5$.

As mentioned in Example~\ref{ax:5.28}, no exact expression is known for any of
the zeros of this polynomial, but numerical techniques show that the zeros of
this polynomial are approximately $-1.67$, $0.51$, $1.40$, $-0.12+1.59i$,
$-0.12-1.59i$.

The software that produces these approximations is accurate to more than three
digits. Thus these approximations are good enough to show that the five numbers
above are distinct. The minimal polynomial of $T$ equals the fifth degree monic
polynomial with these zeros. Now \ref{ax:5.62} shows that $T$ is
diagonalizable.
\end{example}

\begin{example}[Showing that an operator is not diagonalizable]\label{ax:5.61}
Define $T\in\Lin(\F^3)$ by
\[ T(z_1,z_2,z_3)=(6z_1+3z_2+4z_3,6z_2+2z_3,7z_3). \]
The matrix of $T$ with respect to the standard basis of $\F^3$ is
\[ \begin{pmatrix} 6&3&4\\ 0&6&2\\ 0&0&7 \end{pmatrix}. \]
The matrix above is an upper-triangular matrix but is not a diagonal matrix.
Might $T$ have a diagonal matrix with respect to some other basis of $\F^3$?

To answer this question, we will find the minimal polynomial of $T$. First note
that the eigenvalues of $T$ are the diagonal entries of the matrix above (by
\ref{ax:5.41}). Thus the zeros of the minimal polynomial of $T$ are $6,7$
[by \ref{ax:5.27}(a)]. The diagonal of the matrix above tells us that
$(T-6I)^2(T-7I)=0$ (by \ref{ax:5.40}). The minimal polynomial of $T$ has degree
at most $3$ (by \ref{ax:5.22}). Putting all this together, we see that the
minimal polynomial of $T$ is either $(z-6)(z-7)$ or $(z-6)^2(z-7)$.

A simple computation shows that $(T-6I)(T-7I)\ne0$. Thus the minimal polynomial
of $T$ is $(z-6)^2(z-7)$.

Now \ref{ax:5.62} shows that $T$ is not diagonalizable.
\end{example}

\begin{result}[Necessary and sufficient condition for diagonalizability]%
\label{ax:5.62}
Suppose $V$ is finite-dimensional and $T\in\Lin(V)$. Then $T$ is diagonalizable
if and only if the minimal polynomial of $T$ equals
$(z-\lambda_1)\cdots(z-\lambda_m)$ for some list of distinct numbers
$\lambda_1,\dots,\lambda_m\in\F$.
\end{result}

\begin{proof}
First suppose $T$ is diagonalizable. Thus there is a basis $v_1,\dots,v_n$ of
$V$ consisting of eigenvectors of $T$. Let $\lambda_1,\dots,\lambda_m$ be the
distinct eigenvalues of $T$. Then for each $v_j$, there exists $\lambda_k$ with
$(T-\lambda_kI)v_j=0$. Thus
\[ (T-\lambda_1I)\cdots(T-\lambda_mI)v_j=0, \]
which implies that the minimal polynomial of $T$ equals
$(z-\lambda_1)\cdots(z-\lambda_m)$.

To prove the implication in the other direction, now suppose the minimal
polynomial of $T$ equals $(z-\lambda_1)\cdots(z-\lambda_m)$ for some list of
distinct numbers $\lambda_1,\dots,\lambda_m\in\F$. Thus
\[ (T-\lambda_1I)\cdots(T-\lambda_mI)=0. \axtag{ax:5.63} \]

We will prove that $T$ is diagonalizable by induction on $m$. To get started,
suppose $m=1$. Then $T-\lambda_1I=0$, which means that $T$ is a scalar multiple
of the identity operator, which implies that $T$ is diagonalizable.

Now suppose that $m>1$ and the desired result holds for all smaller values of
$m$. The subspace $\range(T-\lambda_mI)$ is invariant under $T$ [this is a
special case of \ref{ax:5.18} with $p(z)=z-\lambda_m$]. Thus $T$ restricted to
$\range(T-\lambda_mI)$ is an operator on $\range(T-\lambda_mI)$.

If $u\in\range(T-\lambda_mI)$, then $u=(T-\lambda_mI)v$ for some $v\in V$, and
\ref{ax:5.63} implies
\[ (T-\lambda_1I)\cdots(T-\lambda_{m-1}I)u=(T-\lambda_1I)\cdots(T-\lambda_mI)v=0.
   \axtag{ax:5.64} \]
Hence $(z-\lambda_1)\cdots(z-\lambda_{m-1})$ is a polynomial multiple of the
minimal polynomial of $T$ restricted to $\range(T-\lambda_mI)$
[by \ref{ax:5.29}]. Thus by our induction hypothesis, there is a basis of
$\range(T-\lambda_mI)$ consisting of eigenvectors of $T$.

Suppose that $u\in\range(T-\lambda_mI)\cap\nul(T-\lambda_mI)$. Then
$Tu=\lambda_mu$. Now \ref{ax:5.64} implies that
\begin{align*}
0 &= (T-\lambda_1I)\cdots(T-\lambda_{m-1}I)u\\
  &= (\lambda_m-\lambda_1)\cdots(\lambda_m-\lambda_{m-1})u.
\end{align*}
Because $\lambda_1,\dots,\lambda_m$ are distinct, the equation above implies that
$u=0$. Hence $\range(T-\lambda_mI)\cap\nul(T-\lambda_mI)=\{0\}$.

Thus $\range(T-\lambda_mI)+\nul(T-\lambda_mI)$ is a direct sum (by
\ref{ax:1.46}) whose dimension is $\dim V$ (by \ref{ax:3.94} and \ref{ax:3.21}).
Hence $\range(T-\lambda_mI)\oplus\nul(T-\lambda_mI)=V$. Every nonzero vector in
$\nul(T-\lambda_mI)$ is an eigenvector of $T$ with eigenvalue $\lambda_m$.
Earlier in this proof we saw that there is a basis of $\range(T-\lambda_mI)$
consisting of eigenvectors of $T$. Adjoining to that basis a basis of
$\nul(T-\lambda_mI)$ gives a basis of $V$ consisting of eigenvectors of $T$. The
matrix of $T$ with respect to this basis is a diagonal matrix, as desired.
\end{proof}

No formula exists for the zeros of polynomials of degree $5$ or greater.
However, the previous result can be used to determine whether an operator on a
complex vector space is diagonalizable without even finding approximations of the
zeros of the minimal polynomial---see Exercise~15.

The next result will be a key tool when we prove a result about the simultaneous
diagonalization of two operators; see \ref{ax:5.76}. Note how the use of a
characterization of diagonalizable operators in terms of the minimal polynomial
(see \ref{ax:5.62}) leads to a short proof of the next result.

\begin{result}[Restriction of diagonalizable operator to invariant subspace]%
\label{ax:5.65}
Suppose $T\in\Lin(V)$ is diagonalizable and $U$ is a subspace of $V$ that is
invariant under $T$. Then $T|_U$ is a diagonalizable operator on $U$.
\end{result}

\begin{proof}
Because the operator $T$ is diagonalizable, the minimal polynomial of $T$ equals
$(z-\lambda_1)\cdots(z-\lambda_m)$ for some list of distinct numbers
$\lambda_1,\dots,\lambda_m\in\F$ (by \ref{ax:5.62}). The minimal polynomial of
$T$ is a polynomial multiple of the minimal polynomial of $T|_U$ (by
\ref{ax:5.31}). Hence the minimal polynomial of $T|_U$ has the form required by
\ref{ax:5.62}, which shows that $T|_U$ is diagonalizable.
\end{proof}

\subsection{Gershgorin Disk Theorem}

\begin{definition}[Gershgorin disks]\label{ax:5.66}
Suppose $T\in\Lin(V)$ and $v_1,\dots,v_n$ is a basis of $V$. Let $A$ denote the
matrix of $T$ with respect to this basis. A \emph{Gershgorin disk} of $T$ with
respect to the basis $v_1,\dots,v_n$ is a set of the form
\[ \biggl\{z\in\F : \abs{z-A_{j,j}}\le\sum_{\substack{k=1\\k\ne j}}^n
   \abs{A_{j,k}}\biggr\}, \]
where $j\in\{1,\dots,n\}$.
\end{definition}

Because there are $n$ choices for $j$ in the definition above, $T$ has $n$
Gershgorin disks. If $\F=\C$, then for each $j\in\{1,\dots,n\}$, the
corresponding Gershgorin disk is a closed disk in $\C$ centered at $A_{j,j}$,
which is the $j^{\text{th}}$ entry on the diagonal of $A$. The radius of this
closed disk is the sum of the absolute values of the entries in row $j$ of $A$,
excluding the diagonal entry. If $\F=\R$, then the Gershgorin disks are closed
intervals in $\R$.

In the special case that the square matrix $A$ above is a diagonal matrix, each
Gershgorin disk consists of a single point that is a diagonal entry of $A$ (and
each eigenvalue of $T$ is one of those points, as required by the next result).
One consequence of our next result is that if the nondiagonal entries of $A$ are
small, then each eigenvalue of $T$ is near a diagonal entry of $A$.

\begin{result}[Gershgorin disk theorem]\label{ax:5.67}
Suppose $T\in\Lin(V)$ and $v_1,\dots,v_n$ is a basis of $V$. Then each
eigenvalue of $T$ is contained in some Gershgorin disk of $T$ with respect to the
basis $v_1,\dots,v_n$.
\end{result}

\begin{proof}
Suppose $\lambda\in\F$ is an eigenvalue of $T$. Let $w\in V$ be a corresponding
eigenvector. There exist $c_1,\dots,c_n\in\F$ such that
\[ w=c_1v_1+\cdots+c_nv_n. \axtag{ax:5.68} \]

Let $A$ denote the matrix of $T$ with respect to the basis $v_1,\dots,v_n$.
Applying $T$ to both sides of the equation above gives
\begin{align*}
\lambda w &= \sum_{k=1}^n c_kTv_k \axtag{ax:5.69}\\
  &= \sum_{k=1}^n c_k\sum_{j=1}^n A_{j,k}v_j\\
  &= \sum_{j=1}^n\Bigl(\sum_{k=1}^n A_{j,k}c_k\Bigr)v_j. \axtag{ax:5.70}
\end{align*}

Let $j\in\{1,\dots,n\}$ be such that
\[ \abs{c_j}=\max\{\abs{c_1},\dots,\abs{c_n}\}. \]
Using \ref{ax:5.68}, we see that the coefficient of $v_j$ on the left side of
\ref{ax:5.69} equals $\lambda c_j$, which must equal the coefficient of $v_j$ on
the right side of \ref{ax:5.70}. In other words,
\[ \lambda c_j=\sum_{k=1}^n A_{j,k}c_k. \]
Subtract $A_{j,j}c_j$ from each side of the equation above and then divide both
sides by $c_j$ to get
\begin{align*}
\abs{\lambda-A_{j,j}} &= \biggl\lvert\sum_{\substack{k=1\\k\ne j}}^n
     A_{j,k}\frac{c_k}{c_j}\biggr\rvert\\
  &\le \sum_{\substack{k=1\\k\ne j}}^n \abs{A_{j,k}}.
\end{align*}
Thus $\lambda$ is in the $j^{\text{th}}$ Gershgorin disk with respect to the
basis $v_1,\dots,v_n$.
\end{proof}

\marginnote{The Gershgorin disk theorem is named for Semyon Aronovich
Gershgorin, who published this result in 1931.}%
Exercise~22 gives a nice application of the Gershgorin disk theorem.

Exercise~23 states that the radius of each Gershgorin disk could be changed to
the sum of the absolute values of corresponding column entries (instead of row
entries), excluding the diagonal entry, and the theorem above would still hold.

% ------------------------------------------------------------------ exercises
\exercisesheading{5D}

\begin{exercise}
Suppose $V$ is a finite-dimensional complex vector space and $T\in\Lin(V)$.
\begin{parts}
\item Prove that if $T^4=I$, then $T$ is diagonalizable.
\item Prove that if $T^4=T$, then $T$ is diagonalizable.
\item Give an example of an operator $T\in\Lin(\C^2)$ such that $T^4=T^2$ and
  $T$ is not diagonalizable.
\end{parts}
\end{exercise}

\begin{exercise}
Suppose $T\in\Lin(V)$ has a diagonal matrix $A$ with respect to some basis of
$V$. Prove that if $\lambda\in\F$, then $\lambda$ appears on the diagonal of $A$
precisely $\dim E(\lambda,T)$ times.
\end{exercise}

\begin{exercise}
Suppose $V$ is finite-dimensional and $T\in\Lin(V)$. Prove that if the operator
$T$ is diagonalizable, then $V=\nul T\oplus\range T$.
\end{exercise}

\begin{exercise}
Suppose $V$ is finite-dimensional and $T\in\Lin(V)$. Prove that the following
are equivalent.
\begin{parts}
\item $V=\nul T\oplus\range T$.
\item $V=\nul T+\range T$.
\item $\nul T\cap\range T=\{0\}$.
\end{parts}
\end{exercise}

\begin{exercise}
Suppose $V$ is a finite-dimensional complex vector space and $T\in\Lin(V)$.
Prove that $T$ is diagonalizable if and only if
\[ V=\nul(T-\lambda I)\oplus\range(T-\lambda I) \]
for every $\lambda\in\C$.
\end{exercise}

\begin{exercise}
Suppose $T\in\Lin(\F^5)$ and $\dim E(8,T)=4$. Prove that $T-2I$ or $T-6I$ is
invertible.
\end{exercise}

\begin{exercise}
Suppose $T\in\Lin(V)$ is invertible. Prove that
\[ E(\lambda,T)=E\Bigl(\tfrac{1}{\lambda},T^{-1}\Bigr) \]
for every $\lambda\in\F$ with $\lambda\ne0$.
\end{exercise}

\begin{exercise}
Suppose $V$ is finite-dimensional and $T\in\Lin(V)$. Let
$\lambda_1,\dots,\lambda_m$ denote the distinct nonzero eigenvalues of $T$.
Prove that
\[ \dim E(\lambda_1,T)+\cdots+\dim E(\lambda_m,T)\le\dim\range T. \]
\end{exercise}

\begin{exercise}
Suppose $R,T\in\Lin(\F^3)$ each have $2,6,7$ as eigenvalues. Prove that there
exists an invertible operator $S\in\Lin(\F^3)$ such that $R=S^{-1}TS$.
\end{exercise}

\begin{exercise}
Find $R,T\in\Lin(\F^4)$ such that $R$ and $T$ each have $2,6,7$ as eigenvalues,
$R$ and $T$ have no other eigenvalues, and there does not exist an invertible
operator $S\in\Lin(\F^4)$ such that $R=S^{-1}TS$.
\end{exercise}

\begin{exercise}
Find $T\in\Lin(\C^3)$ such that $6$ and $7$ are eigenvalues of $T$ and such that
$T$ does not have a diagonal matrix with respect to any basis of $\C^3$.
\end{exercise}

\begin{exercise}
Suppose $T\in\Lin(\C^3)$ is such that $6$ and $7$ are eigenvalues of $T$.
Furthermore, suppose $T$ does not have a diagonal matrix with respect to any
basis of $\C^3$. Prove that there exists $(z_1,z_2,z_3)\in\C^3$ such that
\[ T(z_1,z_2,z_3)=(6+8z_1,7+8z_2,13+8z_3). \]
\end{exercise}

\begin{exercise}
Suppose $A$ is a diagonal matrix with distinct entries on the diagonal and $B$ is
a matrix of the same size as $A$. Show that $AB=BA$ if and only if $B$ is a
diagonal matrix.
\end{exercise}

\begin{exercise}\leavevmode
\begin{parts}
\item Give an example of a finite-dimensional complex vector space and an
  operator $T$ on that vector space such that $T^2$ is diagonalizable but $T$ is
  not diagonalizable.
\item Suppose $\F=\C$, $k$ is a positive integer, and $T\in\Lin(V)$ is
  invertible. Prove that $T$ is diagonalizable if and only if $T^k$ is
  diagonalizable.
\end{parts}
\end{exercise}

\begin{exercise}
Suppose $V$ is a finite-dimensional complex vector space, $T\in\Lin(V)$, and $p$
is the minimal polynomial of $T$. Prove that the following are equivalent.
\begin{parts}
\item $T$ is diagonalizable.
\item There does not exist $\lambda\in\C$ such that $p$ is a polynomial multiple
  of $(z-\lambda)^2$.
\item $p$ and its derivative $p'$ have no zeros in common.
\item The greatest common divisor of $p$ and $p'$ is the constant polynomial $1$.
\end{parts}
\exnote{The \textbf{greatest common divisor} of $p$ and $p'$ is the monic
polynomial $q$ of largest degree such that $p$ and $p'$ are both polynomial
multiples of $q$. The Euclidean algorithm for polynomials (look it up) can
quickly determine the greatest common divisor of two polynomials, without
requiring any information about the zeros of the polynomials. Thus the
equivalence of (a) and (d) above shows that we can determine whether $T$ is
diagonalizable without knowing anything about the zeros of $p$.}
\end{exercise}

\begin{exercise}
Suppose that $T\in\Lin(V)$ is diagonalizable. Let $\lambda_1,\dots,\lambda_m$
denote the distinct eigenvalues of $T$. Prove that a subspace $U$ of $V$ is
invariant under $T$ if and only if there exist subspaces $U_1,\dots,U_m$ of $V$
such that $U_k\subseteq E(\lambda_k,T)$ for each $k$ and
$U=U_1\oplus\cdots\oplus U_m$.
\end{exercise}

\begin{exercise}
Suppose $V$ is finite-dimensional. Prove that $\Lin(V)$ has a basis consisting of
diagonalizable operators.
\end{exercise}

\begin{exercise}
Suppose that $T\in\Lin(V)$ is diagonalizable and $U$ is a subspace of $V$ that
is invariant under $T$. Prove that the quotient operator $T/U$ is a
diagonalizable operator on $V/U$.
\exnote{The quotient operator $T/U$ was defined in Exercise~38 in Section~5A.}
\end{exercise}

\begin{exercise}
Prove or give a counterexample: If $T\in\Lin(V)$ and there exists a subspace $U$
of $V$ that is invariant under $T$ such that $T|_U$ and $T/U$ are both
diagonalizable, then $T$ is diagonalizable.
\exnote{See Exercise~13 in Section~5C for an analogous statement about
upper-triangular matrices.}
\end{exercise}

\begin{exercise}
Suppose $V$ is finite-dimensional and $T\in\Lin(V)$. Prove that $T$ is
diagonalizable if and only if the dual operator $T'$ is diagonalizable.
\end{exercise}

\begin{exercise}
The \emph{Fibonacci sequence} $F_0,F_1,F_2,\dots$ is defined by
\[ F_0=0,\ F_1=1,\ \text{and } F_n=F_{n-2}+F_{n-1} \text{ for } n\ge2. \]
Define $T\in\Lin(\R^2)$ by $T(x,y)=(y,x+y)$.
\begin{parts}
\item Show that $T^n(0,1)=(F_n,F_{n+1})$ for each nonnegative integer $n$.
\item Find the eigenvalues of $T$.
\item Find a basis of $\R^2$ consisting of eigenvectors of $T$.
\item Use the solution to (c) to compute $T^n(0,1)$. Conclude that
  \[ F_n=1/\sqrt5\biggl[\biggl(\frac{1+\sqrt5}{2}\biggr)^{\!n}
     -\biggl(\frac{1-\sqrt5}{2}\biggr)^{\!n}\biggr] \]
  for each nonnegative integer $n$.
\item Use (d) to conclude that if $n$ is a nonnegative integer, then the
  Fibonacci number $F_n$ is the integer that is closest to
  \[ 1/\sqrt5\biggl(\frac{1+\sqrt5}{2}\biggr)^{\!n}. \]
\end{parts}
\exnote{Each $F_n$ is a nonnegative integer, even though the right side of the
formula in (d) does not look like an integer. The number
\[ \frac{1+\sqrt5}{2} \]
is called the \textbf{golden ratio}.}
\end{exercise}

\begin{exercise}
Suppose $T\in\Lin(V)$ and $A$ is an $n$-by-$n$ matrix that is the matrix of $T$
with respect to some basis of $V$. Prove that if
\[ \abs{A_{j,j}}>\sum_{\substack{k=1\\k\ne j}}^n \abs{A_{j,k}} \]
for each $j\in\{1,\dots,n\}$, then $T$ is invertible.
\exnote{This exercise states that if the diagonal entries of the matrix of $T$
are large compared to the nondiagonal entries, then $T$ is invertible.}
\end{exercise}

\begin{exercise}
Suppose the definition of the Gershgorin disks is changed so that the radius of
the $k^{\text{th}}$ disk is the sum of the absolute values of the entries in
column (instead of row) $k$ of $A$, excluding the diagonal entry. Show that the
Gershgorin disk theorem (\ref{ax:5.67}) still holds with this changed definition.
\end{exercise}
